%%%%%%%%%%%%%%%%%%%%%%%%%%%% Document Setup %%%%%%%%%%%%%%%%%%
\documentclass[10pt]{article}
\usepackage[T2A,T1]{fontenc}
\usepackage[utf8]{inputenc}
\usepackage[russian,english]{babel}
\usepackage{CJKutf8}
\AtBeginDvi{\input{zhwinfonts}}
%\usepackage{polyglossia}
\makeatletter
\newlength{\bibhang}
\setlength{\bibhang}{1em}
\newlength{\bibsep}
 {\@listi \global\bibsep\itemsep \global\advance\bibsep by\parsep}
%\newenvironment{bibsection}%
 %       {\vspace{-\baselineskip}\begin{list}{}{%
   %    \setlength{\leftmargin}{\bibhang}%
     %  \setlength{\itemindent}{-\leftmargin}%
       %\setlength{\itemsep}{\bibsep}%
       %\setlength{\parsep}{\z@}%
        %\setlength{\partopsep}{0pt}%
        %\setlength{\topsep}{0pt}}}
        %{\end{list}\vspace{-.6\baselineskip}}
\makeatother
\reversemarginpar
\usepackage{geometry}\geometry{a4paper,total={170mm,257mm},left=1.18in,right=1.18in,top=0.79in,bottom=0.60in}
%\usepackage[paper=letterpaper,
           %\includefoot, % Uncomment to put page number above margin
   %         marginparwidth=0.6in, % Length of section titles
         %   marginparsep=0.05in, % Space between titles and text
      %      margin=1.18in, % 1 inch margins
            %includemp]{geometry}
%\usepackage{eurosans}
\setlength{\parindent}{0in}
\usepackage{enumerate} 
\usepackage{etaremune}
\usepackage{calc}
\usepackage{bbding}
\usepackage{paralist}
\usepackage{fancyhdr,lastpage}
\pagestyle{fancy}
%\pagestyle{empty} % Uncomment this to get rid of page numbers
\fancyhf{}\renewcommand{\headrulewidth}{0pt}
\fancyfootoffset{\marginparsep+\marginparwidth}
\newlength{\footpageshift}
\setlength{\footpageshift}
          {0.5\textwidth+1.\marginparsep+0.5\marginparwidth-2in}
\lfoot{\hspace{\footpageshift}%
       \parbox{4in}{\, \hfill %
                    \arabic{page} of \protect\pageref*{LastPage} % +LP
% \arabic{page} % -LP
                    \hfill \,}}
\usepackage{color,hyperref}
\usepackage{graphics}
\usepackage{wrapfig}
\usepackage{longtable}
\usepackage[usenames,dvipsnames]{xcolor}
\definecolor{darkblue}{rgb}{0.0,0.0,0.3}
\definecolor{mygreen}{RGB}{44,204,148}
\definecolor{darkgreen}{RGB}{13,107,27}
\definecolor{lightgreen}{RGB}{0,170,37}
\definecolor{provinces}{RGB}{11,100,132}
\definecolor{mypurple}{RGB}{77,24,137}
\definecolor{haiyan}{RGB}{8,150,178}
\definecolor{anticred}{RGB}{165,62,11}
\hypersetup{colorlinks,breaklinks,
            linkcolor=darkblue,urlcolor=darkblue,
            anchorcolor=darkblue,citecolor=darkblue}
%%%%%%%%%%%%%%%%%%%%%%%% End Document Setup %%%%%%%%%%%%%%%%%%%
%%%%%%%%%%%%%%%%%%%%%%%%%%% Helper Commands %%%%%%%%%

\newcommand{\makeheading}[1]%
        {\hspace*{-\marginparsep minus \marginparwidth}%
         \begin{minipage}[t]{\textwidth+\marginparwidth+\marginparsep}%
                {\large \bfseries #1}\\[-0.15\baselineskip]%
                 \rule{\columnwidth}{1pt}%
         \end{minipage}}
\renewcommand{\section}[2]%
        {\pagebreak[2]\vspace{1.3\baselineskip}%
         \phantomsection\addcontentsline{toc}{section}{#1}%
         \hspace{0in}%
         \marginpar{
         \raggedright \scshape #1}#2}
% An itemize-style list with lots of space between items
\newenvironment{outerlist}[1][\enskip\textbullet]%
        {\begin{itemize}[#1]}{\end{itemize}%
         \vspace{-.6\baselineskip}}
% An environment IDENTICAL to outerlist that has better pre-list spacing
% when used as the first thing in a \section
\newenvironment{lonelist}[1][\enskip\textbullet]%
        {\vspace{-\baselineskip}\begin{list}{#1}{%
        \setlength{\partopsep}{0pt}%
        \setlength{\topsep}{0pt}}}
        {\end{list}\vspace{-.6\baselineskip}}

% An itemize-style list with little space between items
\newenvironment{innerlist}[1][\enskip\textbullet]%
        {\begin{compactitem}[#1]}{\end{compactitem}}
% An environment IDENTICAL to innerlist that has better pre-list spacing
% when used as the first thing in a \section
\newenvironment{loneinnerlist}[1][\enskip\textbullet]%
        {\vspace{-\baselineskip}\begin{compactitem}[#1]}
        {\end{compactitem}\vspace{-.6\baselineskip}}
% To add some paragraph space between lines.
% This also tells LaTeX to preferably break a page on one of these gaps
% if there is a needed pagebreak nearby.
\newcommand{\blankline}{\quad\pagebreak[2]}
% Uses hyperref to link DOI
\newcommand\doilink[1]{\href{http://dx.doi.org/#1}{#1}}
\newcommand\doi[1]{doi:\doilink{#1}}
% For \url{SOME_URL}, links SOME_URL to the url SOME_URL
\providecommand*\url[1]{\href{#1}{#1}}
% Same as above, but pretty-prints SOME_URL in teletype fixed-width font
\renewcommand*\url[1]{\href{#1}{\texttt{#1}}}

\include{tikz-3dplot}
\usepackage{tikz,tikz-3dplot}
\usepackage{amssymb}
\usepackage{xifthen}
\usepackage{pstricks,pst-solides3d}
\usetikzlibrary{chains}
\usepackage{pgfplots}
\pgfplotsset{
compat=1.16,
    compat/path replacement=1.5.1,
}
\usepackage{pgfmath}
\usetikzlibrary{matrix}
\usetikzlibrary{arrows,decorations.pathmorphing,backgrounds,positioning,fit,petri}
\usepgfmodule{decorations}
\usepgflibrary{decorations.shapes} 
\usetikzlibrary[decorations.shapes] 
\usetikzlibrary{mindmap}
\usepgflibrary{patterns}
\usetikzlibrary[patterns]
\usetikzlibrary{petri}
\usepgflibrary{plothandlers}
\usetikzlibrary{plothandlers}
\usepackage{pgfpages}
\usepgflibrary{arrows}
\usetikzlibrary{arrows}
\usetikzlibrary{trees}
\usetikzlibrary{topaths}
\usetikzlibrary{3d}
\usetikzlibrary{positioning}
\usetikzlibrary{fit}
\usepgflibrary{shapes.misc}
\usetikzlibrary{shapes.misc}
\usepgflibrary{shadings}
\usetikzlibrary{shadings}
\usepgfmodule{shapes}
\usepgfplotslibrary{external}
\usetikzlibrary{pgfplots.external}
\usepgfplotslibrary{colorbrewer}
\usepgfplotslibrary{patchplots}
\usetikzlibrary{plotmarks}
\usepackage{color} 
\usepackage{xcolor} 
%\pgfpagesuselayout{resize to}[a4paper,border shrink=20mm]
\pgfplotsset{colormap={cool}{rgb255(0cm)=(255,255,255); rgb255(1cm)=(0,128,255); rgb255(2cm)=(255,0,255)}}
  \pgfplotsset{colormap={violet}{rgb255=(25,25,122) color=(white) rgb255=(238,140,238)}}
  \pgfplotsset{  colormap={bluered}{rgb255(0cm)=(0,0,180); rgb255(1cm)=(0,255,255); rgb255(2cm)=(100,255,0);rgb255(3cm)=(255,255,0); rgb255(4cm)=(255,0,0); rgb255(5cm)=(128,0,0)}}
 \pgfplotsset{/pgfplots/colormap={winter}{rgb255=(0,0,255) rgb255=(0,255,128)}}
 \pgfplotsset{/pgfplots/colormap={spring}{rgb255=(255,0,255) rgb255=(255,255,0)}}
 \pgfplotsset{colormap={CM}{rgb=(0,0,1) color=(black) rgb255=(238,140,238)}}
 
\usetikzlibrary{spy}
 \pgfplotsset{every patch/.style={miter limit=50},}
\usepgfplotslibrary{colormaps}


\newcommand*\subfigheight{\linewidth}
\newcommand*\subfigwidth{\linewidth}


\begin{document}
\selectlanguage{russian}

\begin{tikzpicture} \begin{axis}[small,colorbar sampled line,title=Experiment №5802,xlabel={\tiny{Time, h}},
    ylabel={\tiny{Immersion depth, mm}}, zlabel={\tiny{$C_{eq}$, MPa}}, xlabel style={sloped like x axis},
    ylabel style={sloped},colormap name=bluered,grid=major,view={210}{30},minor x tick num=10,minor y tick num=10]
    \addplot3+ [mesh,scatter,
        samples=48,domain=0:1, ] coordinates {
(0,0.25,0.38) (20,0.52,0.20) (40	,0.58,0.19)  (60,0.60,0.18) (80,0.61,0.17) (100,0.62,0.16) (120,0.63,0.15) (140,0.64,0.14)
(0.5,0.4,0.38) (1.5,0.45,0.20) (2.0,0.46,0.19) (2.5,0.47,0.18) (3.0,0.49,0.17) (3.5,0.50,0.16) (4,0.55,0.15) (5,0.60,0.14) 
(0,0.0,0.25) (10,0.15,0.23) (20,0.20,0.22) (30,0.30,0.21) (40,0.40,0.20) (50,0.50,0.19) (55,0.60,0.18) (60,0.70,0.17)
(0,0.0,0.35) (2.0,0.04	,0.31) (4,0.08,0.29) (5,0.10,0.27) (6,0.12,0.25) (7,0.14,0.23) (8,0.15,0.21) (10,0.16,0.20)
(0,0.0,0.34) (2.0,0.10,0.31) (3.0,0.20,0.28) (4,0.30,0.26) (5,0.35	,0.24) (6,0.4,0.23) (7,0.45,0.22) (8,0.50,0.21)
(0,0.0,0.33) (5,0.20,0.30) (7,0.30,0.29) (11	,0.35,0.28) (13,0.40,0.26) (15,0.45,0.24) (17,0.47,0.22) (19,0.52,0.20)  
};

\end{axis}
\end{tikzpicture}
	\hskip 5pt % insert a non-breaking space of specified width.	
\begin{tikzpicture} \begin{axis}[small=4cm,colorbar sampled line,title=Experiment №5803,xlabel={\tiny{Time, h}},
    ylabel={\tiny{Immersion depth, mm}}, zlabel={\tiny{($C_{eq}$), MPa}}, xlabel style={sloped like x axis},
    ylabel style={sloped like y axis}, zticklabel style={rotate=90,
    	/pgf/number format/precision=3, 
    	/pgf/number format/fixed,
        /pgf/number format/fixed zerofill,
},yticklabel style={sloped like y axis},colormap/PuOr,grid=major,view={210}{30},minor x tick num=10,minor y tick num=10]
    \addplot3+ [mesh,scatter,
        samples=42,domain=0:1, ] coordinates {
(0,0.6,0.110) (20,0.9,0.070) (40,1.0,0.062) (60,1.2,0.060) (80,1.3,0.059) (100,1.4,0.058) (120,1.5,0.057) 
(0,0.50,0.110) (2,0.65,0.105) (4,0.70,0.100) (6,0.73,0.080) (8,0.77,0.075) (12,0.80,0.073) (16,0.85,0.070)
(0,0.40,0.110) (2,0.52,0.100) (4,0.62,0.090) (6,0.65,0.088) (8,0.70,0.085) (12,0.73,0.083) (16,0.75,0.080) 
(0,0.50,0.130) (2,0.58,0.110) (3,0.63,0.106) (5,0.65,0.103) (7,0.70,0.100) (10,0.72,0.095) (12,0.75,0.090) 
(0,0.40,0.170) (2,0.50,0.155) (4,0.55,0.150) (6,0.58,0.135) (8,0.62,0.120) (12,0.67,0.110) (16,0.70,0.100) 
(0,0.51,0.133) (2,0.59,0.115) (3,0.64,0.104) (5,0.67,0.096) (7,0.72,0.090) (10,0.74,0.083) (12,0.77,0.080)
};

\end{axis}
\end{tikzpicture}
	\hskip 10pt % insert a non-breaking space of specified width.
	
\begin{tikzpicture} [spy using outlines={red,circle,magnification=1.5, size=2.5cm},connect spies]
\begin{axis}[colorbar sampled line,small,title=Experiment №5804,xlabel={\tiny{Time, h}},
    ylabel={\tiny{Immersion depth, mm}}, zlabel={\tiny{($C_{eq}$), MPa}}, xlabel style={sloped like x axis},
    ylabel style={sloped},colormap/RdGy-9,grid=major,view={210}{30},minor x tick num=10,minor y tick num=10]
    \addplot3+ [mesh,scatter,
        samples=42,domain=0:1, ] coordinates {
(0,0.5,0.09) (50,1.3,0.04) (100,1.5,0.03) (150,1.6,0.025) (200,1.7,0.020) (250,1.8,0.150) (300,1.9,0.100) 
(0,0.8,0.8) (2,0.92,0.75) (3	,0.95,0.73) (4,0.98,0.65) (5,1.0,0.62) (6,1.05,0.61) (8,1.1,0.59) 
(0,0.55,0.13) (10,0.70,0.10) (20,0.75,0.09) (30,0.80,0.083) (40,0.85,0.077) (50,0.95,0.075) (70,1.00,0.07) 
(0,0.4,0.13) (2,0.5,0.12) (4	,0.55,0.17) (6,0.57,0.14) (8,0.61,0.10) (10,0.65,0.90) (12,0.70,0.083)
(0,0.4,0.140) (2,0.5,0.130) (4,0.52,0.122) (6,0.57,0.115) (8,0.63	,0.080) (10,0.68,0.09) (12,0.70,0.080) 
(0,0.4,0.130) (2,0.5,0.128) (4,0.52,0.120) (6,0.57,0.111) (8,0.63,0.079) (10,0.68,0.086) (12,0.70	,0.077) 
 };
        	\coordinate (spypoint) at (50,0.6,0.16);
        \coordinate (magnifyglass) at (280,1.0,0.5);
\end{axis}
	\spy on (spypoint) in node at (magnifyglass);

\end{tikzpicture}
	\hskip 10pt % insert a non-breaking space of specified width.
	\begin{tikzpicture} \begin{axis}[colorbar sampled line,small,title=Experiment №5805,xlabel={\tiny{Time, h}},
    ylabel={\tiny{Immersion depth, mm}}, zlabel={\tiny{($C_{eq}$) , MPa}}, xlabel style={sloped like x axis},
    ylabel style={sloped},colormap name=hot,grid=major,view={210}{30},minor x tick num=10,minor y tick num=10]
    \addplot3+ [mesh,scatter,
        samples=42,domain=0:1, ] coordinates {
 (0,0.25,0.530) (10,0.38,0.280) (20,0.40,0.260) (30,0.41,0.250) (40,0.45,0.245) (50,0.475,0.241) (60,0.491,0.235) (70,0.495,0.230) 
 (0,0.20,0.45) (2,0.30	,0.32) (4,0.32,0.32) (7,0.34,0.31) (10,0.36,0.300) (11,0.38,0.290) (13,0.40,0.280) (15	,0.42,0.270)
 (0,0.25,0.38) (2,0.30,0.31) (4,0.32,0.30) (7,0.34,0.290) (10,0.36,0.280) (11,0.38,0.27) (13,0.40,0.260) (15	,0.42	,0.250) 
 (0,0.20,0.400) (2,0.22,0.30) (3,0.24,0.280) (4,0.26,0.270) (5,0.28,0.260) (6,0.30,0.255) (7	,0.33,0.240) (8,0.35,0.230)
 (0,0.21,0.403) (2,0.23,0.302) (3,0.24,0.283) (4,0.27,0.272) (5,0.29,0.257) (6,0.31,0.253) (7,0.34,0.220) (8,0.36	,0.180) 
 (0,0.26,0.400) (2,0.31,0.33) (3,0.33,0.31) (4,0.35,0.280) (5,0.37,0.27) (6,0.39,0.260) (7,0.40,0.250) (8,0.41,0.240)
};
\end{axis}
\end{tikzpicture}\\
	\hskip 10pt % insert a non-breaking space of specified width.
	\begin{tikzpicture} \begin{axis}[colorbar sampled line,small,title=Experiment №5810,xlabel={\tiny{Time, h}},
    ylabel={\tiny{Immersion depth, mm}}, zlabel={\tiny{($C_{eq}$), MPa}},xlabel style={sloped like x axis},
    ylabel style={sloped},colormap name=winter,grid=major,view={210}{30},minor x tick num=10,minor y tick num=10]
    \addplot3+ [mesh,scatter,
        samples=42,domain=0:1, ] coordinates {
(0,0.5,0.130) (20,1.1,0.060) (40,1.2,0.055) (60,1.22,0.050) (80,1.24,0.049) (100,1.26,0.043) (120,1.28,0.042) (140,1.29,0.041) 
(0,0.3,0.90) (10,0.45,0.85) (20,0.60,0.75) (25,0.62,0.65) (30,0.63,0.55) (40,0.65,0.45) (50,0.68,0.35) (62,0.7,0.20)   
(0,0.15,0.270) (2,0.25,0.220) (4,0.35,0.170) (6,0.45,0.150) (8,0.50,0.130) (10,0.52,0.120) (14,0.53,0.115) (18,0.55,0.112)
(0,0.0,0.160) (1,0.20,0.135) (2,0.25,0.125) (3,0.30,0.115) (4,0.35,0.105) (5,0.55	,0.102) (7,0.65,0.100) (10,0.75,0.095)
(0,0.0,0.170) (1,0.20,0.160) (2,0.25,0.155) (3,0.33,0.150) (4,0.38,0.145) (5,0.41,0.140) (7,0.45,0.135) (8,0.5,0.130)
(0,0.0,0.060) (2,0.15,0.58) (4,0.25,0.43) (5,0.35	,0.42) (7,0.45,0.40) (11,0.60,0.35) (15,0.70,0.30) (18,0.75,0.017)     
};
\end{axis}
\end{tikzpicture}
	\hskip 10pt % insert a non-breaking space of specified width.
	\begin{tikzpicture} \begin{axis}[colorbar sampled line,small,title=Experiment №5812,xlabel={\tiny{Time, h}},
    ylabel={\tiny{Immersion depth, mm}}, zlabel={\tiny{($C_{eq}$) , MPa}},xlabel style={sloped like x axis},
    ylabel style={sloped},colormap/BrBG-9,grid=major,view={210}{30},minor x tick num=10,minor y tick num=10]
    \addplot3+ [mesh,scatter,
        samples=42,domain=0:1, ] coordinates {
(0,0.30,0.300) (10,0.55,0.185) (20,0.60,0.180) (30,0.61,0.170) (40,0.62,0.160) (50,0.63,0.150) (60,0.64,0.100)    
(0,0.0,0.300) (1.5,0.10,0.240) (1.5,0.20,0.220) (2,0.30,0.200) (3	,0.40,0.190) (4,0.50,0.185) (5,0.60,0.180)    
(0,0.50,0.600) (2,0.55,0.420) (4	,0.60,0.350) (6,0.62,0.330) (8,0.64,0.310) (12,0.68,0.290) (15,0.70,0.280)    
(0,0.40,0.250) (2,0.45,0.220) (4,0.50,0.210) (6,0.52,0.200) (8,0.54,0.195) (12,0.58,0.190) (15,0.60,0.180)    
(0,0.30,0.300) (2,0.32,0.280) (3,0.34,0.260) (4,0.36,0.240) (5,0.38,0.230) (6,0.40,0.220) (7,0.42,0.210)    
(0,0.150,0.150) (2,0.130,0.140) (4,0.110,0.130) (6,0.090,0.125) (8,0.070,0.120) (12,0.060,0.115) (15,0.050,0.110)    
};
\end{axis}
\end{tikzpicture}
\\
%	\hskip 10pt % insert a non-breaking space of specified width.
	\begin{tikzpicture} \begin{axis}[colorbar sampled line,small,title=Experiment №5831,xlabel={\tiny{Time, h}},
    ylabel={\tiny{Immersion depth, mm}}, zlabel={\tiny{($C_{eq}$), MPa}},xlabel style={sloped like x axis},
    ylabel style={sloped},colormap/PiYG-11,grid=major,view={240}{30},minor x tick num=10,minor y tick num=10]
    \addplot3+ [mesh,scatter,
        samples=66,domain=0:1, ] coordinates {
(0,0.20,0.33) (10,0.36,0.200) (20,0.38,0.190) (30,0.390,0.180) (40,0.395,0.177) (50,0.397,0.175) (60,0.399,0.170) (70,0.405,0.165) (80,0.407,0.160) (90,0.408,0.155) (100,0.410,0.150)           
(0,0.50,0.152) (1,0.53,0.140) (2,0.57,0.130) (3,0.60,0.120) (4,0.65,0.113) (6,0.67,0.104) (8,0.70,0.098) (10,0.71,0.096) (12,0.75,0.083) (14,0.77,0.075) (16,0.80,0.065) 
(0,0.70,0.100) (2,0.80,0.090) (3,0.90,0.085) (4.5,1.00,0.080) (5.5,1.10,0.075) (5.5,1.12,0.070) (6,1.15,0.065) (6.5,1.2,0.060) (7.0,1.3,0.055) (7.5,1.4,0.052) (8,1.5,0.050) 
(0,0.4,0.150) (2,0.5,0.140) (3,0.55,0.130) (4.5,0.58,0.120) (5.5,0.59,0.110) (5.5,0.6,0.100) (6,0.65,0.098) (6.5,0.7,0.096) (7.0,0.72,0.094) (7.5,0.74,0.092) (8,0.8,0.090) 
(0,0.50,0.150) (1,0.60,0.138) (2,0.65,0.128) (3,0.70,0.118) (4,0.75,0.111) (6,0.77,0.102) (8	,0.80,0.099) (10,0.81,0.097) (12,0.85,0.080) (14,0.87,0.070) (16,0.90,0.060) 
(0,0.30,0.250) (1,0.35,0.230) (2,0.405,0.192) (3,0.45,0.169) (4,0.47,0.145) (6,0.49,0.130) (8,0.51,0.120) (10,0.53,0.109) (12,0.55,0.103) (14,0.57,0.110) (16,0.60,0.100)  
};
\end{axis}
\end{tikzpicture}
	\hskip 10pt % insert a non-breaking space of specified width.
	\begin{tikzpicture} [spy using outlines={red,circle,magnification=1.5, size=2.5cm},connect spies]
	\begin{axis}[colorbar sampled line,small,title=Experiment №5833,xlabel={\tiny{Time, h}},
    ylabel={\tiny{Immersion depth, mm}}, zlabel={\tiny{($C_{eq}$) , MPa}},xlabel style={sloped like x axis},
    ylabel style={sloped},colormap/RdYlBu-4,grid=major,view={240}{30},minor x tick num=10,minor y tick num=10]
    \addplot3+ [mesh,scatter,
        samples=42,domain=0:1, ] coordinates {
 (0,1.0,0.100) (20,1.60,0.080) (40,1.70,0.060) (60,1.75,0.030) (80,1.80,0.020) (100,1.81,0.010)   
 (0,0.5,0.25) (2,0.55,0.20) (6,0.57,0.18) (8,0.62,0.17) (12,0.66,0.16) (18,0.7,0.15) 
 (0,0.300,0.30) (2,0.350,0.26) (4,0.400,0.25) (6,0.440,0.24) (8,0.490,0.22) (10,0.550,0.20) 
 (0,0.31,0.32) (2,0.33	,0.27) (4,0.41,0.24) (6,0.43,0.22) (8,0.47,0.20) (10,0.53,0.18) 
 (0,0.51,0.24) (2,0.54,0.22) (4,0.57,0.21) (6,0.58,0.18) (8,0.60,0.17) (10,0.62,0.16) 
 (0,0.2,0.23) (2,0.350	,0.21) (4,0.47,0.20) (6,0.57,0.19) (8,0.58,0.18) (10,0.61,0.17)  
};
        	\coordinate (spypoint) at (50,0.16,0.16);
        \coordinate (magnifyglass) at (90,1.0,0.16);
\end{axis}
	\spy on (spypoint) in node at (magnifyglass);
\end{tikzpicture}

%	\begin{flushleft}
	%	\selectlanguage{russian}\textbf{
%		\textbf{Fig. 1. Time series analysis of the variations of the long-term equivalent cohesion of soils ($C_{eq}$) under increasing depth of the samples. Series-II (№1-8). 3D: \LaTeX{}}
%	\end{flushleft}
\end{document}
%Рис. 1. Временные ряды изменения велиhины длительного эквивалентного сцепления грунтов ($C_{eq}$) с увелиhением глубины погружения образцов. Серии испытаний-I (№1-8). 3D-визуализация: \LaTeX{}}\\\vspace{6pt}
